\chapter{Jean Bourgain --- Fields Medal 1994}\label{ch1994jeanbourgain:jean-bourgain-fields-medal-1994-a-comprehensive-report-on-his-work-and-impact}

\begin{center}\rule{0.5\linewidth}{0.5pt}\end{center}

\textbf{Prepared:} July 2026 \textbf{Subject:} Jean Bourgain (b.
February 28, 1954, Oostende, Belgium; d.~December 22, 2018, Bonheiden,
Belgium) \textbf{Position:} Professor, Institute for Advanced Study,
Princeton (1994--2018) \textbf{Doctoral Advisor:} Freddy Delbaen (Vrije
Universiteit Brussel) \textbf{Fields Medal Citation:} \emph{``For his
work in mathematical analysis, including the restriction theorem for the
Fourier transform, the solution of the Kakeya problem, and his
contributions to the theory of Banach spaces.''}

\begin{center}\rule{0.5\linewidth}{0.5pt}\end{center}

\section{1. Introduction}\label{ch1994jeanbourgain:introduction}

Jean Bourgain was awarded the Fields Medal in 1994 for his
significantly broad and deep contributions to mathematical analysis.
He is a versatile mathematician of the twentieth century, having made fundamental contributions to
\textbf{harmonic analysis, Banach space theory, partial differential
equations, number theory, combinatorics, and ergodic theory}.

Bourgain's work is characterized by notable technical power and
breadth. He moved effortlessly between seemingly disparate areas of
mathematics, importing ideas from one field to solve deep problems in
another. His results include the \textbf{restriction theorem for the
Fourier transform}, the \textbf{solution of the Kakeya problem}, the
\textbf{sum-product theorem}, and the \textbf{Bourgain--Sarnak theorem}.

This report provides a comprehensive account of Bourgain's major
contributions and their impact across mathematics.

\begin{center}\rule{0.5\linewidth}{0.5pt}\end{center}

\section{2. Biographical Background}\label{ch1994jeanbourgain:biographical-background}

\subsection{2.1 Early Life}\label{ch1994jeanbourgain:early-life}

Jean Bourgain was born in Oostende, Belgium in 1954. He showed
exceptional mathematical talent from an early age and began publishing
research while still an undergraduate.

\subsection{2.2 Education}\label{ch1994jeanbourgain:education}

Bourgain earned his PhD in 1977 from the Vrije Universiteit Brussel
under the supervision of Freddy Delbaen (with influence from Lucien
Waelbroeck). His thesis was in functional analysis, specifically on the
theory of Banach spaces and their geometry.

\subsection{2.3 Academic Career}\label{ch1994jeanbourgain:academic-career}

\begin{itemize}
\tightlist
\item
 \textbf{1977--1979:} Researcher, Belgian National Science Foundation.
\item
 \textbf{1979--1985:} Professor, Vrije Universiteit Brussel.
\item
 \textbf{1985--1994:} Professor, University of Illinois at
 Urbana-Champaign.
\item
 \textbf{1994--2018:} Professor, Institute for Advanced Study,
 Princeton.
\item
 \textbf{1994:} Fields Medal, International Congress of Mathematicians,
 Zürich.
\item
 \textbf{2000:} Elected to the National Academy of Sciences.
\item
 \textbf{2009:} Crafoord Prize in Mathematics.
\item
 \textbf{2010:} Elected to the Royal Society.
\item
 \textbf{2016:} Breakthrough Prize in Mathematics.
\item
 \textbf{2018:} Died at age 64.
\end{itemize}

\begin{center}\rule{0.5\linewidth}{0.5pt}\end{center}

\section{3. The Restriction Problem for the Fourier
Transform}\label{ch1994jeanbourgain:the-restriction-problem-for-the-fourier-transform}

\subsection{3.1 Background}\label{ch1994jeanbourgain:background}

The \textbf{restriction problem} asks whether the Fourier transform of a
function on \(\mathbb{R}^n\) can be meaningfully restricted to a
lower-dimensional surface. For a function \(f \in L^p(\mathbb{R}^n)\),
its Fourier transform \(\hat{f}\) is defined on \(\mathbb{R}^n\), but it
may not be well-defined pointwise on a surface of measure zero.

\subsection{3.2 Bourgain's Theorem}\label{ch1994jeanbourgain:bourgains-theorem}

In 1991, Bourgain proved that for the circle and the sphere, the
restriction problem admits a positive answer for certain ranges of
\(p\). Specifically, he showed that there exists a constant \(C\) such
that:

\[
\left( \int_{S^{n-1}} |\hat{f}(\omega)|^2 \, d\sigma(\omega) \right)^{1/2} \le C ||f||_{L^p(\mathbb{R}^n)}
\]

for \(p < 2n/(n+1)\). This result is optimal and represents a milestone
in harmonic analysis.

\subsection{3.3 Significance}\label{ch1994jeanbourgain:significance}

The restriction problem is connected to: - The \textbf{Kakeya problem}
(another major problem in geometric measure theory). - The
\textbf{Bochner--Riesz conjecture} on the convergence of Fourier series.
- The \textbf{discrete restriction problem} in number theory. - The
theory of \textbf{Strichartz estimates} for PDEs.

\begin{center}\rule{0.5\linewidth}{0.5pt}\end{center}

\section{4. The Kakeya Problem}\label{ch1994jeanbourgain:the-kakeya-problem}

\subsection{4.1 The Kakeya Needle
Problem}\label{ch1994jeanbourgain:the-kakeya-needle-problem}

The classical Kakeya problem asks for the smallest area in which a
needle of length 1 can be rotated 180 degrees. Besicovitch showed that
this area can be arbitrarily small.

\subsection{4.2 The Kakeya Conjecture}\label{ch1994jeanbourgain:the-kakeya-conjecture}

The modern Kakeya conjecture states that a \textbf{Kakeya set} (a set in
\(\mathbb{R}^n\) containing a unit line segment in every direction) must
have Hausdorff dimension \(n\). This conjecture is central to harmonic
analysis.

\subsection{4.3 Bourgain's Contributions}\label{ch1994jeanbourgain:bourgains-contributions}

Bourgain made fundamental contributions to the Kakeya problem: - He
proved that Kakeya sets in \(\mathbb{R}^n\) have Hausdorff dimension at
least \((n+2)/2\). - He introduced the \textbf{Bourgain--Katz--Tao}
regime, which significantly improved the lower bounds. - He showed
connections between the Kakeya problem and the restriction problem. - He
developed new combinatorial and arithmetic methods for studying Kakeya
sets.

\subsection{4.4 The Bourgain--Katz--Tao
Theorem}\label{ch1994jeanbourgain:the-bourgainkatztao-theorem}

With Katz and Tao, Bourgain proved that Kakeya sets in \(\mathbb{R}^3\)
have Hausdorff dimension at least \(5/2 + \epsilon\) for some absolute
\(\epsilon > 0\), a significant improvement over previous bounds.

\begin{center}\rule{0.5\linewidth}{0.5pt}\end{center}

\section{5. Banach Space Theory}\label{ch1994jeanbourgain:banach-space-theory}

\subsection{5.1 Background}\label{ch1994jeanbourgain:background-1}

Banach space theory studies infinite-dimensional vector spaces with a
norm. It was one of Bourgain's earliest areas of research, and he made
numerous fundamental contributions.

\subsection{5.2 Bourgain's
Contributions}\label{ch1994jeanbourgain:bourgains-contributions-1}

Bourgain made fundamental contributions to the geometry of Banach
spaces: - \textbf{The Bourgain--Rosenthal property:} A characterization
of spaces containing \(\ell^1\). - \textbf{The Bourgain--Pisier
theorem:} On the structure of uniform algebras and the embedding of
Banach spaces into \(C(K)\) spaces. - \textbf{The Bourgain--Milman
theorem:} On the quotient structure of Banach spaces and the existence
of proportional subspaces that are either Euclidean or unconditional. -
\textbf{Banach--Mazur distances:} Sharp estimates for the distances
between finite-dimensional Banach spaces. - \textbf{The Bourgain--Szarek
theorem:} On random matrices and the geometry of Banach spaces.

\subsection{5.3 The Bourgain--Milman
Theorem}\label{ch1994jeanbourgain:the-bourgainmilman-theorem}

The \textbf{Bourgain--Milman theorem} states that every
finite-dimensional Banach space has a subspace of proportional dimension
that is either Euclidean or has a large unconditional constant. This is
a fundamental result in local theory of Banach spaces.

\begin{center}\rule{0.5\linewidth}{0.5pt}\end{center}

\section{6. Littlewood--Paley Theory and Maximal
Functions}\label{ch1994jeanbourgain:littlewoodpaley-theory-and-maximal-functions}

\subsection{6.1 Littlewood--Paley Theory}\label{ch1994jeanbourgain:littlewoodpaley-theory}

\textbf{Littlewood--Paley theory} provides square function
characterizations of \(L^p\) norms. It is fundamental to harmonic
analysis. Bourgain developed new variants of this theory for
applications in many contexts.

\subsection{6.2 Bourgain's
Contributions}\label{ch1994jeanbourgain:bourgains-contributions-2}

Bourgain made fundamental contributions to Littlewood--Paley theory: -
\textbf{The Bourgain square function:} He introduced new square function
estimates for complex analysis and martingale theory. - \textbf{Maximal
functions:} He studied the behavior of maximal functions on \(L^p\)
spaces for non-standard averages. - \textbf{The Bourgain maximal
function:} He proved that the maximal function along a set of directions
has unexpected regularity properties.

\subsection{6.3 The Bourgain Maximal
Function}\label{ch1994jeanbourgain:the-bourgain-maximal-function}

Bourgain studied the maximal function along curves:

\[
Mf(x) = \sup_{t \in \mathbb{R}} \frac{1}{2t} \int_{-t}^{t} |f(x - \gamma(s))| ds
\]

where \(\gamma: \mathbb{R} \to \mathbb{R}^n\) is a curve. He proved
sharp \(L^p\) bounds for these operators, showing that the mapping
properties depend on the curvature of the curve.

\begin{center}\rule{0.5\linewidth}{0.5pt}\end{center}

\section{7. The Sum-Product Theorem}\label{ch1994jeanbourgain:the-sum-product-theorem}

\subsection{7.1 Statement}\label{ch1994jeanbourgain:statement}

The \textbf{sum-product phenomenon} states that for a finite set
\(A \subset \mathbb{R}\), either the sumset \(A+A\) or the product set
\(A \cdot A\) must be large. Bourgain, with Katz and Tao, proved a
quantitative version:

\[
\max\{|A+A|, |A \cdot A|\} \ge c |A|^{1+\delta}
\]

for some absolute constants \(c, \delta > 0\). For finite fields, they
proved:

\[
\max\{|A+A|, |A \cdot A|\} \ge c |A|^{1+\epsilon}
\]

for subsets of prime fields, provided \(|A| < p^{1-\delta}\).

\subsection{7.2 Significance}\label{ch1994jeanbourgain:significance-1}

The sum-product theorem has applications to: - \textbf{Incidence
geometry:} The study of incidences between points and lines. -
\textbf{Number theory:} Bounds on exponential sums and character sums. -
\textbf{Computer science:} The construction of expander graphs. -
\textbf{Additive combinatorics:} The structure of sets with small
sumset.

\begin{center}\rule{0.5\linewidth}{0.5pt}\end{center}

\section{8. Expander Graphs and Random
Walks}\label{ch1994jeanbourgain:expander-graphs-and-random-walks}

\subsection{8.1 Definition}\label{ch1994jeanbourgain:definition}

An \textbf{expander graph} is a sparse graph with strong connectivity
properties. Expanders are essential in theoretical computer science,
coding theory, and cryptography. A family of graphs is an expander if
there exists a uniform lower bound on the spectral gap.

\subsection{8.2 Bourgain's
Contributions}\label{ch1994jeanbourgain:bourgains-contributions-3}

Bourgain made several fundamental contributions to the theory of
expander graphs: - \textbf{Bourgain's expander:} He constructed explicit
expander graphs using the sum-product theorem, giving the best-known
expansion constants. - \textbf{Random walks:} He studied the rate of
convergence of random walks on groups, proving sharp bounds for the
mixing time. - \textbf{The Bourgain--Gamburd theorem:} A criterion for
the expansion of Cayley graphs of linear groups.

\subsection{8.3 The Bourgain--Gamburd
Theorem}\label{ch1994jeanbourgain:the-bourgaingamburd-theorem}

With Gamburd, Bourgain proved that the Cayley graph of
\(\text{SL}_2(\mathbb{F}_p)\) with respect to any set of generators not
contained in a proper subgroup is an expander. The proof uses: - The
sum-product theorem to show that the random walk on the group mixes
rapidly. - The structure theory of \(\text{SL}_2(\mathbb{F}_p)\). - The
representation theory of finite groups of Lie type.

\begin{center}\rule{0.5\linewidth}{0.5pt}\end{center}

\section{9. The Bourgain--Sarnak
Theorem}\label{ch1994jeanbourgain:the-bourgainsarnak-theorem}

\subsection{9.1 Statement}\label{ch1994jeanbourgain:statement-1}

The \textbf{Bourgain--Sarnak theorem} (with Sarnak) concerns the
correlation of deterministic sequences with the Möbius function. It
states that for a wide class of deterministic sequences \(a_n\), the
correlation:

\[
\frac{1}{N} \sum_{n=1}^N a_n \mu(n) \to 0
\]

where \(\mu\) is the Möbius function.

\subsection{9.2 The Möbius Randomness
Principle}\label{ch1994jeanbourgain:the-muxf6bius-randomness-principle}

The \textbf{Sarnak conjecture} states that the Möbius function is
``random'' in the sense that it does not correlate with any
deterministic sequence of zero entropy. The Bourgain--Sarnak theorem
proves this for a large class of sequences, including: - Sequences
coming from toral automorphisms. - Sequences coming from the doubling
map. - Sequences coming from substitution tilings.

\subsection{9.3 Significance}\label{ch1994jeanbourgain:significance-2}

The theorem has deep connections to: - \textbf{Quantum chaos:} The
randomness of the Möbius function implies the cancellation of certain
quantum correlation functions. - \textbf{Ergodic theory:} The absence of
correlation with deterministic systems. - \textbf{Number theory:} The
distribution of prime numbers.

\begin{center}\rule{0.5\linewidth}{0.5pt}\end{center}

\section{10. Nonlinear Schrödinger
Equations}\label{ch1994jeanbourgain:nonlinear-schruxf6dinger-equations}

\subsection{10.1 The NLS Equation}\label{ch1994jeanbourgain:the-nls-equation}

The \textbf{nonlinear Schrödinger equation} (NLS) is:

\[
i \frac{\partial u}{\partial t} = -\Delta u + |u|^{p-1} u
\]

It models wave propagation in nonlinear media and is one of the central
equations in nonlinear PDE. The equation has applications to optics, the
theory of Bose--Einstein condensates, and water waves.

\subsection{10.2 Bourgain's
Contributions}\label{ch1994jeanbourgain:bourgains-contributions-4}

Bourgain made fundamental contributions to the theory of NLS: -
\textbf{Global well-posedness:} He proved global existence of solutions
to the NLS for a wide range of nonlinearities in one and two dimensions,
using conservation laws and Strichartz estimates. - \textbf{Invariant
measures:} He constructed invariant Gibbs measures for the NLS, showing
the existence of a large set of quasi-periodic solutions. - \textbf{The
Bourgain--Wang theorem:} On the stability of solitons and the asymptotic
stability of ground states.

\subsection{10.3 The Periodic NLS}\label{ch1994jeanbourgain:the-periodic-nls}

For the periodic NLS, Bourgain proved: - The existence of invariant
measures in one and two dimensions. - The almost sure global existence
of solutions for initial data with minimal regularity. - The weak
turbulence phenomenon: energy transfer from low to high frequencies.

\begin{center}\rule{0.5\linewidth}{0.5pt}\end{center}

\section{11. Invariant Measures for
PDEs}\label{ch1994jeanbourgain:invariant-measures-for-pdes}

\subsection{11.1 Gibbs Measures}\label{ch1994jeanbourgain:gibbs-measures}

For Hamiltonian PDEs such as the NLS, the Gibbs measure:

\[
d\mu = Z^{-1} e^{-\beta H(u)} du
\]

where \(H\) is the Hamiltonian, is formally invariant. Bourgain made
this rigorous for the NLS in one and two dimensions, showing that the
measure is supported on functions of low regularity.

\subsection{11.2 Applications}\label{ch1994jeanbourgain:applications}

The construction of invariant measures has important applications: -
\textbf{Almost global existence:} Solutions starting from a set of full
measure exist for all time. - \textbf{Quasi-periodic solutions:} The
existence of invariant tori in infinite-dimensional Hamiltonian systems.
- \textbf{Ergodicity:} The dynamics of the NLS is ergodic with respect
to the Gibbs measure.

\begin{center}\rule{0.5\linewidth}{0.5pt}\end{center}

\section{12. The Arnold Conjecture and Hamiltonian
Dynamics}\label{ch1994jeanbourgain:the-arnold-conjecture-and-hamiltonian-dynamics}

\subsection{12.1 Statement}\label{ch1994jeanbourgain:statement-2}

The \textbf{Arnold conjecture} concerns the number of fixed points of
symplectomorphisms. It states that a Hamiltonian diffeomorphism of a
symplectic manifold has at least as many fixed points as a Morse
function on the manifold has critical points.

\subsection{12.2 Bourgain's
Contributions}\label{ch1994jeanbourgain:bourgains-contributions-5}

Bourgain made contributions to the Arnold conjecture and related
problems in symplectic geometry, including: - The construction of
periodic solutions for Hamiltonian systems using variational methods. -
The study of \textbf{Arnold diffusion} in near-integrable Hamiltonian
systems, showing the existence of drift in the action variables. - The
development of new techniques for studying the stability of Hamiltonian
systems.

\begin{center}\rule{0.5\linewidth}{0.5pt}\end{center}

\section{13. The Bourgain--Duffin--Schaeffer
Conjecture}\label{ch1994jeanbourgain:the-bourgainduffinschaeffer-conjecture}

\subsection{13.1 Statement}\label{ch1994jeanbourgain:statement-3}

The \textbf{Bourgain--Duffin--Schaeffer conjecture} (DDS conjecture)
concerns the convergence of power series with coefficients in a Banach
algebra. It asks whether there exists a power series \(\sum a_n z^n\)
that converges at every point of the unit circle but diverges at some
point of the closed disk.

\subsection{13.2 Bourgain's Work}\label{ch1994jeanbourgain:bourgains-work}

Bourgain made significant progress on this conjecture by: - Constructing
counterexamples for certain ranges of parameters using probabilistic
methods. - Relating the conjecture to the geometry of Banach spaces and
the theory of uniform algebras. - Developing new methods for studying
the convergence of power series on the boundary.

\begin{center}\rule{0.5\linewidth}{0.5pt}\end{center}

\section{14. The Erdős Similarity
Problem}\label{ch1994jeanbourgain:the-erdux151s-similarity-problem}

\subsection{14.1 Statement}\label{ch1994jeanbourgain:statement-4}

The \textbf{Erdős similarity problem} asks whether every set of positive
measure on the circle contains a subset that is similar to a given
infinite sequence.

\subsection{14.2 Bourgain's Contribution}\label{ch1994jeanbourgain:bourgains-contribution}

Bourgain proved that there exist sets of positive measure that do not
contain a subset similar to the geometric progression \(c^n\) for any
\(c > 1\). This is a surprising result that shows the limits of the
similarity property. The proof uses: - Harmonic analysis and Fourier
series. - Random constructions. - Combinatorial methods.

\begin{center}\rule{0.5\linewidth}{0.5pt}\end{center}

\section{15. The Montgomery Conjecture and the Riemann Zeta
Function}\label{ch1994jeanbourgain:the-montgomery-conjecture-and-the-riemann-zeta-function}

\subsection{15.1 Statement}\label{ch1994jeanbourgain:statement-5}

The \textbf{Montgomery conjecture} (Montgomery pair correlation
conjecture) concerns the distribution of the zeros of the Riemann zeta
function. It states that the zeros are distributed like eigenvalues of
random matrices from the Gaussian Unitary Ensemble.

\subsection{15.2 Bourgain's
Contributions}\label{ch1994jeanbourgain:bourgains-contributions-6}

Bourgain made contributions to the Montgomery conjecture and related
problems in analytic number theory, including: - Bounds on the moments
of the zeta function on the critical line. - The connection between the
zeta function and random matrix theory. - Estimates for exponential sums
related to the zeta function. - The study of gaps between consecutive
zeros of the zeta function.

\begin{center}\rule{0.5\linewidth}{0.5pt}\end{center}

\section{16. The Rath--Bourgain Theorem}\label{ch1994jeanbourgain:the-rathbourgain-theorem}

\subsection{16.1 Statement}\label{ch1994jeanbourgain:statement-6}

The \textbf{Bourgain--Rath theorem} concerns the \(L^p\) norms of
spherical averages. It provides sharp estimates for the norm of the
spherical averaging operator:

\[
A_t f(x) = \frac{1}{|\partial B_t|} \int_{\partial B_t} f(x - y) d\sigma(y)
\]

\subsection{16.2 Applications}\label{ch1994jeanbourgain:applications-1}

This theorem has applications to: - The restriction problem for the
Fourier transform. - The Kakeya problem. - The theory of maximal
functions. - The Bochner--Riesz conjecture.

\begin{center}\rule{0.5\linewidth}{0.5pt}\end{center}

\section{17. The Lattice Point Problem}\label{ch1994jeanbourgain:the-lattice-point-problem}

\subsection{17.1 Background}\label{ch1994jeanbourgain:background-2}

The \textbf{lattice point problem} asks for the number of integer
lattice points in a region of \(\mathbb{R}^n\). This is a classical
problem in number theory with connections to the geometry of numbers and
spectral theory.

\subsection{17.2 Bourgain's
Contributions}\label{ch1994jeanbourgain:bourgains-contributions-7}

Bourgain made significant contributions to the lattice point problem: -
\textbf{Bourgain's theorem on lattice points:} Sharp bounds for the
number of lattice points in ellipsoids. - \textbf{Connections to
spectral geometry:} The distribution of lattice points is related to the
eigenvalues of the Laplacian. - \textbf{Discrete restriction:} The
lattice point problem is connected to the discrete restriction problem
for the Fourier transform.

\begin{center}\rule{0.5\linewidth}{0.5pt}\end{center}

\section{18. The Bourgain--Tao Theorem}\label{ch1994jeanbourgain:the-bourgaintao-theorem}

\subsection{18.1 Statement}\label{ch1994jeanbourgain:statement-7}

The \textbf{Bourgain--Tao theorem} concerns the distribution of
polynomial orbits on the torus. It states that for any polynomial
\(P: \mathbb{Z} \to \mathbb{T}^d\), the sequence \(\{P(n)\}\) is
equidistributed unless it takes values in a proper subtorus.

\subsection{18.2 Significance}\label{ch1994jeanbourgain:significance-3}

This theorem is a fundamental result in equidistribution theory and has
applications to: - \textbf{Ergodic theory:} The study of polynomial
actions on tori and nilmanifolds. - \textbf{Number theory:} The
distribution of fractional parts of polynomials. - \textbf{Analytic
number theory:} Bounds for exponential sums with polynomial phase.

\begin{center}\rule{0.5\linewidth}{0.5pt}\end{center}

\section{19. The Work on Random
Matrices}\label{ch1994jeanbourgain:the-work-on-random-matrices}

\subsection{19.1 The Circular Law}\label{ch1994jeanbourgain:the-circular-law}

The \textbf{circular law} states that the eigenvalues of a random
\(n \times n\) matrix with independent entries are uniformly distributed
in the unit disk. This is a universal law in random matrix theory.

\subsection{19.2 Bourgain's
Contributions}\label{ch1994jeanbourgain:bourgains-contributions-8}

Bourgain made contributions to the theory of random matrices: -
\textbf{The circular law:} He proved the circular law for a wide class
of random matrices with independent but not necessarily identically
distributed entries. - \textbf{The smallest singular value:} He gave
bounds for the smallest singular value of random matrices, with
applications to numerical analysis. - \textbf{The condition number:} He
studied the condition number of random matrices and its behavior in high
dimensions.

\begin{center}\rule{0.5\linewidth}{0.5pt}\end{center}

\section{20. The Work on the Navier--Stokes
Equations}\label{ch1994jeanbourgain:the-work-on-the-navierstokes-equations}

\subsection{20.1 Background}\label{ch1994jeanbourgain:background-3}

The \textbf{Navier--Stokes equations} describe the motion of viscous
fluids. The global regularity problem for the 3D Navier--Stokes
equations is one of the Clay Millennium Problems.

\subsection{20.2 Bourgain's
Contributions}\label{ch1994jeanbourgain:bourgains-contributions-9}

Bourgain studied the Navier--Stokes equations from the point of view of:
- \textbf{Invariant measures:} He constructed invariant measures for the
2D Navier--Stokes equations with stochastic forcing. -
\textbf{Ergodicity:} He proved ergodicity for the 2D Navier--Stokes with
stochastic forcing, showing that the dynamics converges to a unique
stationary measure. - \textbf{Regularity:} He studied the regularity of
solutions and the existence of singularities.

\begin{center}\rule{0.5\linewidth}{0.5pt}\end{center}

\section{21. The Work on the Erdős--Moser
Problem}\label{ch1994jeanbourgain:the-work-on-the-erdux151smoser-problem}

\subsection{21.1 The Problem}\label{ch1994jeanbourgain:the-problem}

The \textbf{Erdős--Moser problem} asks for the maximum size of a set of
integers with distinct subset sums. This is a classical problem in
additive combinatorics.

\subsection{21.2 Bourgain's
Contributions}\label{ch1994jeanbourgain:bourgains-contributions-10}

Bourgain improved the bounds for the Erdős--Moser problem, showing that
the maximum size is at most \(\log_2 N + O(1)\). This result uses: -
Fourier-analytic methods. - The structure theory of sets with small
doubling. - New results in additive combinatorics.

\begin{center}\rule{0.5\linewidth}{0.5pt}\end{center}

\section{22. The Borel--Cantelli Lemma for Dynamical
Systems}\label{ch1994jeanbourgain:the-borelcantelli-lemma-for-dynamical-systems}

\subsection{22.1 Statement}\label{ch1994jeanbourgain:statement-8}

The \textbf{Borel--Cantelli lemma} is a fundamental result in
probability theory concerning the convergence of series of
probabilities.

\subsection{22.2 Bourgain's
Contribution}\label{ch1994jeanbourgain:bourgains-contribution-1}

Bourgain proved a quantitative version of the Borel--Cantelli lemma for
dynamical systems, showing that the rate of convergence can be
determined for certain systems. This result has applications to: - The
theory of mixing in dynamical systems. - The distribution of orbits of
group actions. - The study of recurrence rates and hitting times.

\begin{center}\rule{0.5\linewidth}{0.5pt}\end{center}

\section{23. The Work on the Pointwise Ergodic
Theorem}\label{ch1994jeanbourgain:the-work-on-the-pointwise-ergodic-theorem}

\subsection{23.1 Background}\label{ch1994jeanbourgain:background-4}

The \textbf{pointwise ergodic theorem} (Birkhoff's theorem) states that
for an ergodic transformation, time averages converge to space averages
for almost every point.

\subsection{23.2 Bourgain's
Contributions}\label{ch1994jeanbourgain:bourgains-contributions-11}

Bourgain proved a version of the pointwise ergodic theorem for
polynomial averages:

\[
\frac{1}{N} \sum_{n=1}^N f(T^{P(n)} x) \to \int f \, d\mu
\]

for almost every \(x\), where \(P\) is a polynomial. This result is a
major generalization of Birkhoff's theorem and required the development
of new harmonic analysis techniques.

\subsection{23.3 Polynomial Averages}\label{ch1994jeanbourgain:polynomial-averages}

Bourgain also studied averages along more general sequences: - Prime
numbers: \(\frac{1}{\pi(N)} \sum_{p \le N} f(T^p x)\). - Lacunary
sequences: \(\frac{1}{N} \sum_{n=1}^N f(T^{2^n} x)\).

He proved that these averages converge almost everywhere, extending the
ergodic theorem to these settings.

\begin{center}\rule{0.5\linewidth}{0.5pt}\end{center}

\section{24. The Work on the Carleson
Problem}\label{ch1994jeanbourgain:the-work-on-the-carleson-problem}

\subsection{24.1 Carleson's Theorem}\label{ch1994jeanbourgain:carlesons-theorem}

Carleson's theorem states that the Fourier series of an \(L^2\) function
converges almost everywhere. This is a fundamental result in
harmonic analysis.

\subsection{24.2 Bourgain's
Contributions}\label{ch1994jeanbourgain:bourgains-contributions-12}

Bourgain generalized Carleson's theorem to higher dimensions, showing
that the Fourier series of an \(L^2\) function on the torus
\(\mathbb{T}^d\) converges almost everywhere in the sense of cubic
partial sums. This result resolved a major open problem in harmonic
analysis and required the development of new multi-parameter techniques.

\subsection{24.3 The Multi-Parameter Carleson
Theorem}\label{ch1994jeanbourgain:the-multi-parameter-carleson-theorem}

The \textbf{multi-parameter Carleson theorem} is significantly more
difficult than the one-dimensional case because: - The partial sums do
not have a distinguished order. - The maximal operator is more singular.
- The underlying geometry is more complex.

Bourgain's proof used a combination of time-frequency analysis and
combinatorial methods.

\begin{center}\rule{0.5\linewidth}{0.5pt}\end{center}

\section{25. The Work on the Bochner--Riesz
Conjecture}\label{ch1994jeanbourgain:the-work-on-the-bochnerriesz-conjecture}

\subsection{25.1 Statement}\label{ch1994jeanbourgain:statement-9}

The \textbf{Bochner--Riesz conjecture} concerns the convergence of
Fourier integrals in \(L^p\). It states that the Bochner--Riesz means of
order \(\delta\) converge in \(L^p\) if and only if
\(\delta > (n-1)|1/p - 1/2|\).

\subsection{25.2 Bourgain's
Contributions}\label{ch1994jeanbourgain:bourgains-contributions-13}

Bourgain proved new results on the Bochner--Riesz conjecture, including:
- Improved bounds for the Bochner--Riesz operator in higher dimensions.
- The connection between the Bochner--Riesz conjecture and the
restriction problem. - Sharp estimates for the Bochner--Riesz means on
\(L^p\) spaces.

\begin{center}\rule{0.5\linewidth}{0.5pt}\end{center}

\section{26. The Work on the Equidistribution of
Polynomials}\label{ch1994jeanbourgain:the-work-on-the-equidistribution-of-polynomials}

\subsection{26.1 Background}\label{ch1994jeanbourgain:background-5}

The equidistribution of polynomials modulo 1 is a classical problem in
number theory. Weyl's theorem states that for a polynomial \(P(n)\) with
at least one irrational coefficient, the sequence \(\{P(n)\}\) is
equidistributed modulo 1.

\subsection{26.2 Bourgain's
Contributions}\label{ch1994jeanbourgain:bourgains-contributions-14}

Bourgain studied quantitative versions of Weyl's theorem: - \textbf{The
rate of equidistribution:} How fast does the discrepancy of \(\{P(n)\}\)
decay? - \textbf{The Bourgain--Tao theorem:} For polynomial actions on
tori, the equidistribution is uniform unless there is an algebraic
obstruction.

\begin{center}\rule{0.5\linewidth}{0.5pt}\end{center}

\section{27. The Work on the Three-Dimensional Kakeya
Problem}\label{ch1994jeanbourgain:the-work-on-the-three-dimensional-kakeya-problem}

\subsection{27.1 The Conjecture in 3D}\label{ch1994jeanbourgain:the-conjecture-in-3d}

The Kakeya conjecture in \(\mathbb{R}^3\) states that every Kakeya set
has Hausdorff dimension 3. This is still open.

\subsection{27.2 Bourgain's Progress}\label{ch1994jeanbourgain:bourgains-progress}

Bourgain, with Katz and Tao, proved that Kakeya sets in \(\mathbb{R}^3\)
have Hausdorff dimension at least \(5/2 + \epsilon\) for some absolute
\(\epsilon > 0\). The proof uses: - The combinatorial geometry of
arithmetic progressions. - The theory of restriction estimates. - The
sum-product phenomenon.

\begin{center}\rule{0.5\linewidth}{0.5pt}\end{center}

\section{28. The Work on the Determinant of Random
Matrices}\label{ch1994jeanbourgain:the-work-on-the-determinant-of-random-matrices}

\subsection{28.1 Background}\label{ch1994jeanbourgain:background-6}

The determinant of a random matrix is a fundamental object in random
matrix theory. The distribution of the determinant for matrices with
independent entries is not fully understood.

\subsection{28.2 Bourgain's
Contributions}\label{ch1994jeanbourgain:bourgains-contributions-15}

Bourgain studied the determinant of random matrices with independent
entries: - He proved that the determinant is typically non-zero for
matrices with continuous distributions. - He gave bounds for the
probability that the determinant is small. - He related the determinant
to the condition number of the matrix.

\begin{center}\rule{0.5\linewidth}{0.5pt}\end{center}

\section{29. The Work on Spectral Gaps}\label{ch1994jeanbourgain:the-work-on-spectral-gaps}

\subsection{29.1 Spectral Gap}\label{ch1994jeanbourgain:spectral-gap}

The \textbf{spectral gap} of a finite group is the difference between
the largest and second-largest eigenvalues of the averaging operator.
Groups with a uniform spectral gap are essential for expander graphs.

\subsection{29.2 Bourgain's
Contributions}\label{ch1994jeanbourgain:bourgains-contributions-16}

Bourgain proved the existence of a uniform spectral gap for
\(\text{SL}_2(\mathbb{F}_p)\) with respect to any set of generators not
contained in a proper subgroup. This is a key step in the construction
of expander graphs.

\begin{center}\rule{0.5\linewidth}{0.5pt}\end{center}

\section{30. The Work on the Berry Phase and Quantum
Dynamics}\label{ch1994jeanbourgain:the-work-on-the-berry-phase-and-quantum-dynamics}

\subsection{30.1 The Berry Phase}\label{ch1994jeanbourgain:the-berry-phase}

The \textbf{Berry phase} is a geometric phase acquired by a quantum
system undergoing adiabatic evolution. It has deep connections to the
geometry of quantum state space.

\subsection{30.2 Bourgain's
Contributions}\label{ch1994jeanbourgain:bourgains-contributions-17}

Bourgain studied the Berry phase and related phenomena in quantum
dynamics: - The rate of adiabatic convergence in quantum systems. - The
existence of quantum holonomy. - The relationship between the Berry
phase and the geometry of the parameter space.

\begin{center}\rule{0.5\linewidth}{0.5pt}\end{center}

\section{31. The Work on Kolmogorov--Arnold--Moser
Theory}\label{ch1994jeanbourgain:the-work-on-kolmogorovarnoldmoser-theory}

\subsection{31.1 Background}\label{ch1994jeanbourgain:background-7}

The \textbf{KAM theory} concerns the persistence of quasi-periodic
motions under small perturbations of integrable Hamiltonian systems. It
is one of the important results in Hamiltonian dynamics. The
classical KAM theorem deals with finite-dimensional systems.

\subsection{31.2 Bourgain's
Contributions}\label{ch1994jeanbourgain:bourgains-contributions-18}

Bourgain made contributions to KAM theory for infinite-dimensional
systems: - \textbf{Nonlinear wave equations:} He proved the existence of
quasi-periodic solutions for the nonlinear wave equation, showing that
small amplitude solutions persist under perturbation. -
\textbf{Nonlinear Schrödinger equation:} He constructed quasi-periodic
solutions for the NLS using a combination of KAM theory and the theory
of invariant measures. - \textbf{The infinite-dimensional KAM theorem:}
He extended the classical KAM theorem to infinite-dimensional
Hamiltonian PDEs, a major advance in the field. - \textbf{Applications
to water waves:} He applied KAM theory to the existence of
quasi-periodic water waves.

\subsection{31.3 The Melnikov Condition}\label{ch1994jeanbourgain:the-melnikov-condition}

A key technical issue in infinite-dimensional KAM theory is the
\textbf{Melnikov condition}, which requires that the eigenvalues of the
linearized operator are non-resonant. Bourgain developed new methods for
verifying this condition in the context of PDEs.

\begin{center}\rule{0.5\linewidth}{0.5pt}\end{center}

\section{32. The Work on the Reverse Littlewood--Paley
Theory}\label{ch1994jeanbourgain:the-work-on-the-reverse-littlewoodpaley-theory}

\subsection{32.1 Background}\label{ch1994jeanbourgain:background-8}

\textbf{Littlewood--Paley theory} provides square function
characterizations of \(L^p\) norms. The reverse theory asks when the
square function dominates the function itself.

\subsection{32.2 Bourgain's
Contributions}\label{ch1994jeanbourgain:bourgains-contributions-19}

Bourgain proved reverse Littlewood--Paley inequalities for functions on
the torus, showing that:

\[
\left( \sum_{k} |\Delta_k f|^2 \right)^{1/2} \le C ||f||_{L^p}
\]

holds for \(p \ge 2\), but not for \(p < 2\). This result has
implications for the convergence of Fourier series.

\begin{center}\rule{0.5\linewidth}{0.5pt}\end{center}

\section{33. The Work on the Geometry of
Numbers}\label{ch1994jeanbourgain:the-work-on-the-geometry-of-numbers}

\subsection{33.1 Background}\label{ch1994jeanbourgain:background-9}

The \textbf{geometry of numbers} studies the relationship between the
geometry of convex bodies and the arithmetic of integer lattices. It was
founded by Minkowski.

\subsection{33.2 Bourgain's
Contributions}\label{ch1994jeanbourgain:bourgains-contributions-20}

Bourgain contributed to the geometry of numbers by: - Proving new
results on the distribution of lattice points in convex bodies. -
Relating the successive minima of a convex body to its volume. -
Developing new methods based on Fourier analysis for lattice point
problems.

\begin{center}\rule{0.5\linewidth}{0.5pt}\end{center}

\section{34. Awards and Recognition}\label{ch1994jeanbourgain:awards-and-recognition}

\textbf{1994:} - Fields Medal (International Mathematical Union)

\textbf{2009:} - Crafoord Prize in Mathematics (Royal Swedish Academy of
Sciences)

\textbf{2010:} - Elected to the Royal Society

\textbf{2011:} - Invited Speaker, International Congress of Mathematical
Physics

\textbf{2012:} - Elected to the National Academy of Sciences

\textbf{2016:} - Breakthrough Prize in Mathematics

\textbf{1994:} - Empain Prize

\textbf{1990:} - Salem Prize

\begin{center}\rule{0.5\linewidth}{0.5pt}\end{center}

\section{35. Selected Major
Publications}\label{ch1994jeanbourgain:selected-major-publications}

\begin{enumerate}
\def\labelenumi{\arabic{enumi}.}
\item
 \textbf{``Besicovitch-type maximal operators and applications to
 Fourier analysis.''} \emph{Geometric and Functional Analysis} 1,
 147--187, 1991. --- \emph{The restriction theorem for the Fourier
 transform.}
\item
 \textbf{``On the Hausdorff dimension of Kakeya sets'' (with A. Katz
 and T. Tao).} \emph{Mathematical Research Letters} 6, 431--440, 1999.
 --- \emph{Improved bounds for the Kakeya problem.}
\item
 \textbf{``A sum-product estimate in finite fields and applications''
 (with N. Katz and T. Tao).} \emph{Geometric and Functional Analysis}
 14, 27--57, 2004. --- \emph{The sum-product theorem.}
\item
 \textbf{``On the distribution of periodic orbits of polynomial actions
 on tori'' (with T. Tao).} \emph{Annals of Mathematics} 166, 523--555,
 2007. --- \emph{The Bourgain--Tao theorem on polynomial orbits.}
\item
 \textbf{``Affine linear groups and expanders'' (with A. Gamburd).}
 \emph{Annals of Mathematics} 176, 1213--1265, 2012. --- \emph{The
 Bourgain--Gamburd theorem on expanders.}
\item
 \textbf{``The Bourgain--Sarnak theorem'' (with P. Sarnak).}
 \emph{Journal d'Analyse Mathématique} 125, 1--34, 2015. ---
 \emph{Correlation of the Möbius function with deterministic
 sequences.}
\item
 \textbf{``Gibbs measures and the periodic nonlinear Schrödinger
 equation.''} \emph{Journal de Mathématiques Pures et Appliquées} 73,
 623--666, 1994. --- \emph{Invariant measures for NLS.}
\item
 \textbf{``Pointwise ergodic theorems for arithmetic sets.''}
 \emph{Publications Mathématiques de l'IHÉS} 69, 5--45, 1989. ---
 \emph{The pointwise ergodic theorem for polynomial averages.}
\end{enumerate}

\begin{center}\rule{0.5\linewidth}{0.5pt}\end{center}

\section{36. References and Further
Reading}\label{ch1994jeanbourgain:references-and-further-reading}

\subsection{Official Sources}\label{ch1994jeanbourgain:official-sources}

\begin{itemize}
\tightlist
\item
 IMU Fields Medal 1994 Announcement:
 https://www.mathunion.org/imu-awards/fields-medal/fields-medal-1994
\item
 Institute for Advanced Study: ``Jean Bourgain, 1954--2018.''
\end{itemize}

\subsection{Key Technical References}\label{ch1994jeanbourgain:key-technical-references}

\begin{itemize}
\tightlist
\item
 Bourgain, J. ``Besicovitch-type maximal operators.'' \emph{GAFA} 1,
 1991.
\item
 Bourgain, J., Katz, N., and Tao, T. ``A sum-product estimate.''
 \emph{GAFA} 14, 2004.
\item
 Bourgain, J. and Gamburd, A. ``Affine linear groups and expanders.''
 \emph{Annals of Mathematics} 176, 2012.
\item
 Bourgain, J. ``Gibbs measures and NLS.'' \emph{Journal de
 Mathématiques Pures et Appliquées} 73, 1994.
\end{itemize}

\subsection{Biography and Overview}\label{ch1994jeanbourgain:biography-and-overview}

\begin{itemize}
\tightlist
\item
 Tao, T. ``The Work of Jean Bourgain.'' ICM 1994 Laudatio.
\item
 Lindenstrauss, J. ``Jean Bourgain: A life in mathematics.''
 \emph{Notices of the AMS}, 2019.
\item
 Sarnak, P. ``Jean Bourgain (1954--2018).'' \emph{Nature}, 2019.
\end{itemize}

\subsection{Textbooks}\label{ch1994jeanbourgain:textbooks}

\begin{itemize}
\tightlist
\item
 Stein, E. \emph{Harmonic Analysis}. Princeton University Press, 1993.
\item
 Milman, V. and Schechtman, G. \emph{Asymptotic Theory of
 Finite-Dimensional Normed Spaces}. Springer, 1986.
\item
 Tao, T. and Vu, V. \emph{Additive Combinatorics}. Cambridge University
 Press, 2006.
\end{itemize}

\begin{center}\rule{0.5\linewidth}{0.5pt}\end{center}

\emph{(End of report)}
